% Generated draft: requires mathematical and source review.
\documentclass[11pt]{article}
\usepackage[margin=25mm]{geometry}
\usepackage{fontspec}
\setmainfont{Latin Modern Roman}
\usepackage{amsmath,amssymb,mathtools,unicode-math}
\setmathfont{Latin Modern Math}
\usepackage{enumitem,needspace,xurl,hyperref}
\setlength{\parindent}{0pt}
\setlength{\parskip}{5pt}
\setlength{\emergencystretch}{3em}
\newcommand{\N}{\mathbb{N}}
\newcommand{\Z}{\mathbb{Z}}
\newcommand{\Q}{\mathbb{Q}}
\newcommand{\R}{\mathbb{R}}
\newcommand{\C}{\mathbb{C}}
\newcommand{\F}{\mathbb{F}}
\newcommand{\A}{\mathbb{A}}
\newcommand{\PP}{\mathbb{P}}
\newcommand{\E}{\mathbb{E}}
\newcommand{\eps}{\varepsilon}
\newcommand{\dd}{\,\mathrm{d}}
\begin{document}
\section*{UnsolvedMath ID: 30000189}
\textbf{Source problem code:} catalogue entry
\subsection{Definitions and mathematical statement}
For each integer \(g\) with \(6\leq g\leq10\), write \(X_g=\mathcal A_g^F\) for the first Voronoi compactification named in the source. In the usual moduli interpretation \(\mathcal A_g\) parametrizes principally polarized abelian varieties of dimension \(g\), and the superscript \(F\) denotes the first Voronoi, or perfect-cone, toroidal compactification. The supplied statement says that the ray
\[
 R_g=\R_{\geq0}[C_1]\subseteq\overline{\operatorname{NE}}(X_g)
\]
defines a flipping contraction. Here \([C_1]\) should be the numerical class of the particular curve denoted \(C_1\) in the report, and \(\overline{\operatorname{NE}}\) denotes the closed cone of effective curve classes modulo numerical equivalence. The curve, its intersection numbers and the contraction map are not described in the record.

The principal request is to construct the predicted flip explicitly, for each of these five genera. In the usual canonical-flip formulation, if the indicated contraction is \(f_g:X_g\to Z_g\), this means exhibiting a normal variety \(X_g^+\), a small projective birational contraction \(f_g^+:X_g^+\to Z_g\), and a birational map \(X_g\dashrightarrow X_g^+\) over \(Z_g\), isomorphic in codimension one, such that \(-K_{X_g}\) is \(f_g\)-ample and \(K_{X_g^+}\) is \(f_g^+\)-ample. Small means that the exceptional locus has no divisorial component. These conditions explain the standard term ``flip'' conditionally on the requisite normality and \(\Q\)-Cartier setup; the record does not verify those hypotheses or specify a log boundary or another divisor in place of \(K\).

The cleaned question additionally asks whether the minimal model program on \(\mathcal A_g^F\) can more generally be carried out explicitly: describe the relevant contractions, flips and resulting models, with the governing divisor and stopping criterion specified. The original statement explicitly requests only the flip, so the broader program is retained as a separate extension rather than merged into that single construction. Its scope in \(g\) is not stated beyond the initial range; no claim for all genera is inserted. An abstract existence theorem or identification of a canonical model alone is not the requested explicit description.

The literature assessment checked 2026-08-21 reports that general existence and canonical-model results do not provide the requested explicit flip in the indicated genera and labels this construction unresolved in the checked literature. This dated assessment and the verification field's reported direct-source restoration are retained without independent confirmation.
\subsection{Short English statement}
For \(g=6,7,8,9,10\), explicitly construct the flip of the stated curve ray on the first Voronoi compactification, then investigate an explicit minimal model program there. The curve class, contraction and precise birational setup are missing, so standard flip conditions are given conditionally rather than asserted for unspecified models.
\subsection{Sources}
\begin{itemize}
\item Local UnsolvedMath registry, permanent ID 30000189, source code catalogue entry, ``Explicit Flips on Voronoi Compactifications,'' including original and extended cleaned questions and the reported direct-source verification.
\item \emph{Komplexe Algebraische Geometrie (2005)}, Oberwolfach Reports, pp.\ 407--463, DOI 10.4171/owr/2005/07: \url{https://doi.org/10.4171/owr/2005/07}.
\item Supplied reference ``Shepherd-Barron, perfect forms and the moduli space of abelian varieties'': \url{https://arxiv.org/abs/math/0502362}.
\item Literature triage checked 2026-08-21; registry record updated 2026-09-16. No external source has been visited or independently checked for this draft.
\end{itemize}
\end{document}
